% Source fingerprint: 5cdd0d81d043c04d350859463f9f5ef021ada45e9c19721fbfac6bf4bd71bd18
% T04: raw TeX cells preserved; no numeric highlighting.
\begin{table}[htbp]
\centering\small\renewcommand{\arraystretch}{1.18}\setlength{\tabcolsep}{4pt}
\caption[三个 Banach 空间原理]{三个 Banach 空间原理。三者都使用 Baire 型完备性，但完备性的位置与算子条件不同，不能只检查某个空间完备。}
\label{b1:tab:visual-t04}
\begin{tabularx}{\linewidth}{@{}>{\raggedright\arraybackslash}p{3.1cm}>{\raggedright\arraybackslash}p{4.7cm}>{\raggedright\arraybackslash}X@{}}
\toprule
原理 & 空间与算子条件 & 结论 \\
\midrule
一致有界原理 & \(X\) 为 Banach；\(T\in\mathcal F\subset\mathcal B(X,Y)\) 对每个 \(x\) 的像一致有界 & \(\sup_{T\in\mathcal F}\|T\|<\infty\)；值域不必完备 \\
开映射定理 & \(X,Y\) 为 Banach；\(T:X\to Y\) 为有界线性满射 & 像的开性；有界双射的逆也有界 \\
闭图像定理 & \(X,Y\) 为 Banach；处处定义的线性 \(T\) 具有闭图像 & \(T\) 有界；在乘积空间上把图像完备性用于投影 \\
\bottomrule
\end{tabularx}
\end{table}
